% ============================================================
% 项目全景说明(说人话版) —— sando-core
% 编译方式:XeLaTeX(Overleaf 上把编译器设为 XeLaTeX 即可)。
% 本机(仿真机)没装中文 LaTeX 环境,请传到 Overleaf 编译。
% 现状截至 2026-07-11(数据主要来自 07-04 ~ 07-08 的实验)。
% ============================================================
\documentclass[11pt]{ctexart}

\usepackage[a4paper, margin=2.5cm]{geometry}
\usepackage{amsmath, amssymb}
\usepackage{booktabs}
\usepackage{longtable}
\usepackage{array}
\usepackage{enumitem}
\usepackage{xcolor}
\usepackage[colorlinks=true, linkcolor=blue!60!black, urlcolor=blue!60!black]{hyperref}

\setlist{itemsep=2pt, topsep=4pt}
\newcommand{\code}[1]{\texttt{#1}}

\title{\textbf{带安全证明的无人机街道飞行系统}\\[4pt]
\large 项目全景说明(直白版)\\
\normalsize sando-core \; · \; 现状截至 2026 年 7 月}
\author{}
\date{2026-07-11}

\begin{document}
\maketitle

\begin{abstract}
本文用最直白的语言,把这个项目\textbf{是什么、怎么工作、证明了什么、测出了什么数字、
还差什么},完整梳理一遍。全文尽量不用行话:文献里叫 ``conformal tube'' 的东西,
本文就叫\textbf{预测误差余量};叫 ``evade'' 的叫\textbf{紧急规避};代码里的各种代号
(如 \code{V11}、\code{YOUNG\_TTL})在正文中都用中文直译,并在附录~\ref{app:glossary}
给出对照表,方便回代码里查。
\end{abstract}

\tableofcontents
\newpage

% ============================================================
\section{这个项目在做什么}

\textbf{场景}:一架小型无人机在有行人、车辆的街道上空低空飞行,要从起点飞到终点。
行人和车会动,而且无人机的传感器是``真实''的:只有 $45^\circ$ 的前向视野、会被遮挡、
会漏检、测量有噪声、也不知道每个目标的身份。

\textbf{目标}:又快又安全地到达。项目的头号考核指标叫\textbf{``干净抵达''},
必须同时满足三个条件:
\begin{enumerate}
  \item 到达终点;
  \item 全程零碰撞;
  \item 全程没有触发过任何一次``紧急规避''(即从头到尾都在有安全证明的状态下飞)。
\end{enumerate}

\textbf{核心思路,一句话}:先\textbf{预测}每个行人/车接下来一秒内会到哪里,
让规划器\textbf{提前绕开这些``未来的位置''};在真正执行之前,再用数学方法对规划出的
轨迹做一次\textbf{严格的安全证明}——证明``接下来这一小段飞行,在每一个时刻都不会
进入任何目标的危险范围''。证明通过才飞;通不过就换更保守的动作再证;实在全都证不过,
才做没有保证的紧急规避(这种情况单独记账,算``不干净'')。

\textbf{和一般避障系统的区别}:一般系统是``规划器尽量绕,撞不撞看运气'';
本系统在规划器外面加了一道\textbf{独立的、可证明的安全关卡},并且给出一个可以写成
定理的总保证(见第~\ref{sec:theorem} 节):
\begin{center}
\fbox{\parbox{0.9\textwidth}{\centering
在统计前提成立时,\textbf{一整次飞行中``在已认证状态下发生碰撞''的概率 $\le 5\%$}。\\
(实测中已认证状态下的碰撞为 0;所有残余碰撞都发生在传感器物理上看不见的方向,
且经``全知对照''证实全知也躲不掉,详见第~\ref{sec:honest} 节。)}}
\end{center}

% ============================================================
\section{系统怎么工作:六步流水线}\label{sec:pipeline}

整个系统每 0.3 秒跑一轮下面的循环(``一拍''):

\begin{center}
感知 $\to$ 跟踪与预测 $\to$ 加上预测误差余量 $\to$ 规划 $\to$ 安全证明 $\to$ 决策执行
\end{center}

\subsection{第 1 步:感知(模拟真实传感器)}

无人机不直接读仿真器的``上帝视角''数据,而是通过一个模拟真实传感器的前端
(代码 \code{perception.py},规格按 Intel D435i 深度相机):
\begin{itemize}
  \item 只有前向 $45^\circ$ 视野锥,锥外一律看不见;
  \item 有硬遮挡:被建筑/其他物体挡住就看不见;
  \item 有漏检和位置噪声;
  \item \textbf{不知道目标身份}:这一帧看到的点和上一帧看到的点是不是同一个人,
        要靠最近邻匹配自己猜。
\end{itemize}
这保证了所有实验结论是在``真实感知''条件下取得的,不是全知作弊。
(评测时也可以显式切到全知模式做对照,用于区分``感知问题''和``算法问题''。)

\subsection{第 2 步:跟踪与预测}

每个被看到的目标分配一个\textbf{卡尔曼滤波器}(代码 \code{kf\_tracker.py}),
用``匀加速''运动模型持续估计它的位置、速度、加速度,并外推出未来约 1 秒内的
预测位置曲线
\[
  c(t) = c_0 + v\,t + \tfrac{1}{2}a\,t^2 .
\]
两个重要的诚实规则(都是实验教训,见第~\ref{sec:wins} 节):
\begin{itemize}
  \item \textbf{刚出现的目标不可全信}:观测次数太少时速度估计不可靠,
        这类``新目标''按最坏情况处理,但\textbf{限时}——连续 2 拍还没确认就丢弃,
        不许它无限膨胀;
  \item \textbf{只有``成熟''目标(连续观测 $\ge 4$ 次)的预测才被下游的精细功能采信}。
\end{itemize}

\subsection{第 3 步:预测误差余量(统计标定)}\label{sec:conformal}

预测一定会错。关键问题是:\textbf{错多少?}我们的做法是——不做任何分布假设,
直接\textbf{用数据量出来}:

\begin{enumerate}
  \item 在标定阶段,跑几百集仿真,收集几万条``预测 vs.\ 事后真实位置''的记录,
        每条算出误差 $e$(预测了 $\Delta$ 秒后,真实位置离预测位置多少米);
  \item 给余量定一个随时间变宽的直线形状
        \[
          \text{余量}(\Delta) = b + v\cdot\Delta
          \qquad(\text{截距 } b \text{ + 斜率 } v),
        \]
        行人和车辆分开标,``新目标''和``成熟目标''分开标(误差规律完全不同);
  \item 在标定数据上选出最小的缩放系数,使得 $95\%$ 的记录误差都被余量盖住。
\end{enumerate}

\textbf{它保证什么}:只要新场景和标定场景``统计上同类''
(专业名词叫可交换,直白说就是同一个世界、同一套传感器规律,只是换了随机种子),
那么新场景里\textbf{预测误差有 $95\%$ 的把握不超过余量}。这是一个不依赖任何
分布假设的保证。

\textbf{它怎么用}:第 4、5 步里,每个目标的``危险范围''半径都按预测时长加上这个余量,
即
\[
  R(\Delta) = \underbrace{r_{\text{目标}} + r_{\text{机体}} + d_{\text{安全}}}_{\text{确定性几何}}
            + \underbrace{b + v\cdot\Delta}_{\text{预测误差余量}} .
\]
预测越远,余量越大,``危险范围''越宽——这是诚实的代价。

\textbf{标定方法迭代了三代}(细节见附录~\ref{app:glossary}):
一代是全体记录混在一起取一个分位数(粗);二代按类别、成熟度、预测时长分组细标
(准,但组太多后每组数据不够,出过一系列统计事故);三代(\textbf{现役})把
``形状''(截距和斜率的相对比例)冻结下来,只用新数据估一个总缩放系数——
估计误差只会让余量变宽、绝不会让覆盖率变假。三代标定让行人余量瘦了 $33\%$,
车辆``新目标''余量截距从 14.46 米(离谱)修到 0.71 米,且在保留测试集上
覆盖率 $0.969$(目标 $0.95$)合格。

\subsection{第 4 步:规划(让规划器绕开``未来'')}

规划器用的是浙江大学开源的 EGO-Planner(已剥离 ROS,编译成独立库,
代码 \code{ego/} + \code{ego\_bridge.py})。我们\textbf{不改规划器内部},
只改喂给它的障碍信息:把每个目标\textbf{未来位置的预测}(带余量)转成占据点云
喂进去,规划器就会自然规划出一条\textbf{绕开``目标即将到达的地方''}的轨迹,
而不是等目标到了跟前再躲。

其中一个有效的技巧:把目标摆在``相对运动最近点''时刻的预测位置上
(即它和无人机距离将会最近的那个时刻),规划器就学会了从预测的空隙中穿过,
干净抵达率 $+4.1$ 个百分点。前提就是上面说的``只信成熟目标''——
拿不准的预测宁可不用,错误的预测比没有预测更危险(实验中真撞过一次)。

\subsection{第 5 步:安全证明(核心部件)}\label{sec:cert}

规划器给出轨迹后,\textbf{独立的证明器}(C++ 实现,
代码 \code{cpp/include/sando\_cpp/bernstein\_cert.hpp})要回答:

\begin{center}
\emph{``未来这段时间内,无人机与每个目标的距离,是否在\textbf{每一个时刻}都大于
危险半径 $R$?''}
\end{center}

\textbf{原理(一段话说清)}:无人机轨迹是多项式,目标预测位置也是多项式,
所以``两者距离的平方减去 $R^2$''也是一个关于时间的多项式 $f(t)$。
把 $f(t)$ 改写成一种数学上等价的系数表示(Bernstein 形式),它有个极好的性质:
\textbf{函数在整个区间上的值,永远不超过最大的那个系数}。于是:
\[
  \text{所有系数} \le 0 \;\Longrightarrow\; f(t)\le 0 \text{ 在整段时间成立}
  \;\Longrightarrow\; \text{整段时间距离} \ge R .
\]
一次代数检查,就覆盖了\textbf{区间内的全部无穷多个时刻}——不是抽查几个时间点
(抽查有漏网风险,这正是很多现有系统的隐患)。如果系数有正的但实际是安全的,
就把时间段对半细分再检查,越分越贴近真实曲线(最多 16 层,通常远用不到)。

\textbf{三个可靠性设计}:
\begin{itemize}
  \item \textbf{连续时间、零采样漏洞}:如上,结论对区间内每个时刻成立;
  \item \textbf{浮点数不说谎}:所有运算用``往不利方向取整''的区间算术
        (每步各往外挪一个最小浮点单位),浮点误差只可能让结论更保守,
        绝不可能把危险误判成安全;
  \item \textbf{只判不改}:证明器只输出``通过/不通过 + 裕度'',绝不修改轨迹。
        职责单一,才谈得上可信。
\end{itemize}

\textbf{速度}:单次检查约 7 微秒,一拍内把几十个候选轨迹对全部目标都证一遍毫无压力。
另外,证明器与具体规划器\textbf{解耦}:同一套数学既能证 EGO 的三次样条,
也能证自研规划器的五次多项式(这是支撑性质,不是卖点主线)。

\subsection{第 6 步:决策与执行}\label{sec:decide}

每 0.3 秒(代码 \code{safety\_layer.py} 的统一决策函数):

\begin{enumerate}
  \item 规划器给出一条奔终点的轨迹 $\to$ 全部目标跑一遍安全证明 $\to$
        \textbf{全过就直接飞},本拍结束;
  \item 证不过 $\to$ 生成一批\textbf{候选动作}:方向(原方向/左偏/右偏/爬升)
        $\times$ 速度档(全速到 0.15 倍``爬行挡''),每个候选都独立跑证明,
        在\textbf{全部证得过}的候选里选\textbf{进度最快}的那个执行;
  \item 有``粘性''机制:上一拍的选择在本拍有优先权,避免每拍换主意来回抖;
        另有``逃逸压力''信号:当最近威胁的净空持续收缩时,提前把逃逸类候选排到
        队伍最前(只改排序偏好,每个候选仍必须通过证明);
  \item 候选全军覆没 $\to$ \textbf{紧急规避}:朝远离威胁的方向逃。
        \textbf{这一下没有任何证明保护},所以单独记账,发生一次这集就不算``干净''。
\end{enumerate}

两个与``时间''有关的关键设计(合计是整个调优战役里单笔收益最大的改动,
干净抵达率 $+13.9$ 个百分点):
\begin{itemize}
  \item \textbf{速度换证明窗}:证明需要覆盖的时间窗 = 一拍时长 + 当前速度下的刹停时间。
        飞得慢 $\to$ 刹停快 $\to$ 窗短 $\to$ 预测误差余量小 $\to$ 更容易证过。
        于是``降速''成了换取``可证明性''的连续杠杆;
  \item \textbf{爬行挡}:候选速度网格里加一个 0.15 倍速的最低档。全速证不过时
        还能慢慢挪,而不是停在原地变成活靶。
\end{itemize}

% ============================================================
\section{总保证:两半拼成一个定理}\label{sec:theorem}

第~\ref{sec:conformal} 节的统计余量和第~\ref{sec:cert} 节的几何证明,
各自只保证一半:
\begin{itemize}
  \item 统计半边:``预测不会错得比余量更多''——$95\%$ 概率;
  \item 几何半边:``只要预测没错出余量,就绝对不撞''——$100\%$ 确定。
\end{itemize}
两半拼起来(文档 \code{docs/composition-theorem.md},2026-07-04 定版):

\begin{center}
\fbox{\parbox{0.92\textwidth}{
\textbf{定理(直白版)}:标定场景与测试场景统计上同类时,对一个新场景,
\[
  P\bigl(\text{全程在已认证状态下飞行时发生碰撞}\bigr) \;\le\; \varepsilon = 5\% .
\]
紧急规避时段不在此保证内(单独计量、单独报告)。}}
\end{center}

\textbf{让定理落地的关键一招——``重新证明的节奏''}:
如果余量要覆盖整整 0.85 秒的预测误差,$95\%$ 分位的余量高达 3.18 米——
无人机在街道上根本没法动。但系统\textbf{每 0.3 秒就重新证明一次},
所以任何碰撞发生前的 0.4 秒内必然刚刚通过过一次证明;于是余量只需覆盖
0.4 秒内的预测误差:\textbf{3.18 米 $\to$ 1.02 米},缩小 3 倍,系统才真正可用。
这一步同时把统计的``可交换单元''从聚在一起的记录修正为整段场景,
保留测试覆盖率从 $0.890$ 提到 $0.987$。

\textbf{三本账,防止``证了一半冒充全证''}:
\begin{enumerate}
  \item \textbf{已认证飞行}:上面的 $\varepsilon$ 保证覆盖;
  \item \textbf{沿用旧轨迹}(重规划偶尔失败时继续飞上一条已证轨迹):
        频率可测量,按测得频率单独摊账;
  \item \textbf{紧急规避}:无保证,按暴露时长如实报告。
        (这是结构性的:凡是能证明安全的动作都不会走到这一步。)
\end{enumerate}

% ============================================================
\section{实验结果}\label{sec:results}

\subsection{先说清楚:三个``测量面''}\label{sec:surfaces}

同一套算法在三种环境下测,\textbf{数字不能互相混用},每个数字都必须标明来自哪个面:

\begin{center}
\begin{tabular}{@{}p{3.2cm}p{6.2cm}p{4.6cm}@{}}
\toprule
测量面 & 是什么 & 用途 \\
\midrule
快速回放面 & 在预先录制的真实群体轨迹上离线重放,几百集/波,几十分钟出结果 & 日常迭代、标定、消融 \\
真引擎面 & MetaUrban 完整物理引擎在线跑 & 行为结论的最终复核 \\
3D 渲染面 & 带画面的完整仿真,出视频 & 给人看的演示与取证 \\
\bottomrule
\end{tabular}
\end{center}

\textbf{已知坑}:3D 渲染面目前只读一代标定文件,二、三代标定和部分新开关在渲染面
\textbf{不生效}——所以渲染面数字与评测面数字不可比,调参前必须先分清作用面。

\subsection{主结果:干净抵达率(快速回放面,每档 345 集)}

\begin{center}
\begin{tabular}{@{}lccl@{}}
\toprule
配置 & 干净抵达 & 紧急规避次数 & 说明 \\
\midrule
战役起点(一代标定) & $51.0\%$ & 1290 & 出发点 \\
\textbf{现役最优栈} & $\mathbf{68.4\%}$ & \textbf{531} & 到达 $95.1\%$,\textbf{0 碰撞} \\
视野内全知(对照) & $78.3\%$ & 260 & 感知不变、预测全对的天花板 \\
$360^\circ$ 全知(对照) & $69.3\%$ & 178 & 见下方说明 \\
\bottomrule
\end{tabular}
\end{center}

三点解读:
\begin{itemize}
  \item 现役最优栈把``多余的紧急规避''压掉了 $59\%$(1290 $\to$ 531),
        同时保持全程零碰撞——这正是``干净抵达''哲学:安全不靠疯狂急刹;
  \item \textbf{天花板以 $78.3\%$ 为准}:它保持同样的 $45^\circ$ 视野、只把预测换成
        完美预测,度量的是``预测和决策还有多少可挖''。剩余差距
        ($68.4\% \to 78.3\%$)主要是成熟行人在 3--4 米迎面几何下诚实余量的代价;
  \item \textbf{有趣的``不知者无罪''效应}:$360^\circ$ 全知反而比视野内全知更多急刹
        ——因为它看得见侧后方来车,会提前躲。这说明``更多信息''不等于``更平顺'',
        是论文的讨论点。
\end{itemize}

\subsection{真引擎面复核}

\begin{itemize}
  \item 全栈在真引擎上干净抵达 $64.3\%$,与快速回放面的 $64.3\%$
        (同口径配置)\textbf{精确互证}——回放面结论可信;
  \item \textbf{20 集验收}:全栈共 7 次碰撞,其中 \textbf{6 次全知对照也在同一集死}
        (出生点旁车碾压、高速横穿等,物理上无解),真正归咎于感知栈的只有 1 次;
        排除物理无解的集后干净抵达 $9/14 = 64.3\%$;
  \item 版本阶梯回归(每版本 10 集、场景序列固定):``新目标限时''改动在真引擎复现了
        回放面的收益(急刹 $-36\%$ vs.\ 回放面 $-29\%$),证明改进不是回放面的
        过拟合产物。
\end{itemize}

\subsection{与原生规划器的对比(核心卖点)}

同样的感知、同样的场景,``我们的完整系统'' vs.\ ``原生 EGO-Planner 直接飞'':

\begin{center}
\begin{tabular}{@{}lcc@{}}
\toprule
指标 & 我们 & 原生 EGO \\
\midrule
碰撞(真实感知+动态,304 集) & $4/304 = 1.3\% \le \varepsilon$ & $18/304 = 5.9\%$ \\
碰撞(全知感知各臂,1216 集总账) & \textbf{全臂 0 碰撞 0 近失} & --- \\
最小碰撞时间余量(反应裕度) & 0.59 s & 0.22 s \\
速度代价 & \multicolumn{2}{c}{比原生慢 $5.7\%$} \\
\bottomrule
\end{tabular}
\end{center}

即:\textbf{碰撞降到 1/4 以下、反应裕度 2.7 倍,代价只有 $5.7\%$ 的速度}。
诚实项:按标准舒适性口径,``猛度''(加加速度)指标原生更优——我们的规避动作
仍比原生猛,这是下一步架构工作的目标之一(见第~\ref{sec:ongoing} 节)。
另有 3D 渲染面的论文级对比:弯道新地图上\textbf{原生撞毁、我们以 0.83 米净空
安全到达}(视频 \code{out/PROOF1\_safety\_curvemap.mp4})。

\subsection{强化学习验证臂(支撑性质)}

把同一道安全关卡套在一个纯强化学习(PPO)策略外面,验证``这道关卡不挑规划器'':
\begin{itemize}
  \item 关卡开/关对比:碰撞 $6.75\%$ vs.\ $14.25\%$(400 集,Fisher 检验
        $p = 0.0007$),且开关卡\textbf{还提高}了到达率;
  \item 训练时把标定冻结进训练环境,最优版本碰撞 $6.0\%$ vs.\ 裸策略 $16.25\%$;
  \item 结论:安全关卡对自研规划器、开源规划器、端到端学习策略\textbf{三类}都有效。
        (此线已归档,只作支撑,不再投入。)
\end{itemize}

\subsection{平顺性战役(姿态)}

问题:急刹时机体猛俯仰(倾角变化率 182°/s vs.\ 原生 122°/s),画面难看且危险。
根因:``悬停保持''的参考点是\textbf{冻结}的,6 m/s 的动量撞上一个不动的参考点
自然猛。修法:把参考点换成``按最大减速度连续后退''的\textbf{刹车参考场}
——物理上同样刹停,但参考轨迹不跳变。效果(3D 渲染面固定靶场):
姿态差距 59 $\to$ 12°/s、净空反而 1.30 $\to$ 1.95 米、还快 2.3 秒
——\textbf{平顺不但免费,还倒贴安全和速度}(因为过冲导致的反复修正消失了)。
配合逃逸压力提前触发,一个靶场上姿态与原生\textbf{完全打平}(144 vs.\ 144°/s)
且净空 3 倍。

% ============================================================
\section{一路上的关键改进(问题 $\to$ 原因 $\to$ 修法 $\to$ 效果)}\label{sec:wins}

这些是调优战役中被数据坐实的结构性胜利,每条都很直白:

\begin{enumerate}
  \item \textbf{别把新数据当旧数据用}。
        决策代码假设观测是 0.3 秒前的(历史遗留),实际感知是当拍的——
        白白多算 0.3 秒的预测误差余量。改成 0.05 秒后:急刹 $-74$ 次,
        干净抵达 $+2.6$ 个百分点。
  \item \textbf{没确认的目标不许无限膨胀}。
        一个只见过一两次的``疑似目标''丢失后还能滑行 8 拍,期间按最坏速度
        膨胀成一个巨大的禁飞圈。改成``2 拍内没确认就删''(多目标跟踪的教科书做法):
        急刹 $-350$ 次。
  \item \textbf{同一份误差不许算两次}。
        新目标的标定误差记录本身已经包含了目标移动,部署时又叠加``最坏速度
        $\times$ 时间''——等于运动误差收了两遍钱。拆掉后车辆新目标的余量截距
        从 14.46 米降到 0.70 米(此前一个环岛场景的禁飞圈直径 22.7 米,荒谬):
        急刹 $-195$ 次。
  \item \textbf{速度换证明窗 + 爬行挡}(见第~\ref{sec:decide} 节):
        干净抵达 $+13.9$ 个百分点,整场战役单笔最大。
  \item \textbf{预测``最近点''占据 + 只信成熟目标}:让规划器穿预测缝隙,
        $+4.1$ 个百分点;全量采信新目标的版本真撞过(出生期速度低估被当真),
        限成熟目标后零碰撞且分数更高。
  \item \textbf{刹车参考场救姿态}(上一节)。
  \item \textbf{垂直逃逸的教训}(失败但重要):被人群包围之后再想爬升逃逸,
        \textbf{永远证不过}——上升过程横向暴露在密集禁飞圈里。结论:想从头顶飞过,
        必须在包围圈闭合\textbf{之前}预防性触发,且必须先证明``降落走廊''也安全。
        方向性机动(横穿的目标从其身后过、同向的从头顶过、迎面的横向躲)
        已列为新决策层的候选生成规则,不做旧架构的补丁。
\end{enumerate}

同样重要的是\textbf{被否决的方案}(各有铁证,防止将来重蹈覆辙):
子拍决策(不满足证明前提,撞过 1 次)、按滤波器方差自适应余量(长预测时方差爆炸)、
车辆提前转正(物理上需要 3 个位移样本)、按``自知不确定度''过滤僵尸目标
(滑行目标的不确定度膨胀恰好使其隐身)等。

% ============================================================
\section{诚实边界(明确承认的局限)}\label{sec:honest}

项目哲学:\textbf{边界是主张的一部分}。物理上躲不掉的就承认躲不掉,
用对照实验把``可修的''和``不可修的''分开,不假装没有边界。

\begin{enumerate}
  \item \textbf{视野外快车 = 物理地板}。残余碰撞几乎全部是 $45^\circ$ 视野锥外、
        8.3 m/s 的侧后方快车(无人机才 3 m/s)。$360^\circ$ 全知对照\textbf{一次都
        没能减少}:1.9 秒预警 $\times$ 3 m/s 机动能力,对 8.3 m/s 的扫掠没有逃逸解。
        信息论审查确认只有两类机制能碰它(视野外证据如影子/第二传感器;
        或检测前就降速)——单传感器设定下这是\textbf{物理边界},写入失效包络,不再追。
  \item \textbf{紧急规避无保证}。这是结构性的(能证明的都不会走到这步),
        按暴露时长如实报告;评测显示规避与碰撞集中在同一批退化场景。
  \item \textbf{保证是``分布内''的}。换场景生成器、换传感器规律,必须重新标定
        (标定协议已冻结存档)。标定必须与部署的节奏、感知分布一致——
        这是早期踩过的最大的坑(标定用 0.3 秒节奏、部署 0.1 秒,余量整体失真)。
  \item \textbf{渲染演示面与评测面不同步}(见第~\ref{sec:surfaces} 节):
        演示视频里的行为不能直接当评测证据;演示专用的``人群不停走''开关
        改变了人群统计规律,\textbf{禁止}用于任何安全声明。
  \item \textbf{环境是平面的}。评测环境里高度自由度受限,使得``马路快车''成为致命项;
        真实 3D 产品可以爬升规避,论文需注脚。
  \item \textbf{舒适性的``猛度''指标原生更优}(已如实报告),
        根治靠新架构而不是继续打补丁。
\end{enumerate}

% ============================================================
\section{正在进行的工作}\label{sec:ongoing}

\subsection{新决策层:直接在``被证明安全的轨迹集合''里规划}

\textbf{动机}:现在的决策层是``离散挑候选''(几十个方向$\times$速度组合逐个证明),
焊在连续的规划器上,接缝处天然会犹豫、抖动——一切平滑技巧都是在给一个近视的
仲裁者打补丁。深度调研(104 个研究代理)也确认:根治 = \textbf{在认证集合内规划},
仲裁降为兜底。

\textbf{方案}(代码 \code{local\_lattice.py},设计文档
\code{docs/design-CPL-v3-cert-local-planner.md}):每拍生成一批短轨迹基元,
每条 = ``先飞一小段 + 随后有限猛度刹停 + 停稳保持'',\textbf{整条(含刹停段)
一次性证明},在全部证得过的集合里按字典序(安全$\to$进度$\to$平顺)选优。
好处:任何时刻系统都握着一条``已被证明能安全停下来''的完整退路
(专业上叫递归可行性)。

\textbf{进展与代价}:
\begin{itemize}
  \item 可靠性已锁死:$10^4$ 场景\textbf{零假证明};
  \item 但发现一笔沉重的``随时能刹停税'':要求整条刹停段都安全,意味着行人
        1.75 米内连悬停都证不过(旧决策层可以 0.10 米贴身而过,17 倍差距,
        已在受控几何实验里量化到米):到达率从 8/8 掉到 1/8;
  \item 三个设计级修法方向(已诊断指向,未实施):
        ① 刹停窗内不按``全向最坏速度''膨胀目标(消除双重保守);
        ② 把``完全刹停''弱化为``减速到可重新认证的状态'';
        ③ 主动提前飞越(证得过的前进候选几乎都是爬升类,印证方向性机动的价值)。
\end{itemize}
\textbf{当前裁决}:新决策层暂不替换现役栈;渲染面与评测面现役仍是旧决策层最优栈。

\subsection{第四代标定(已完成,待命)}

直接在保留数据上优化\textbf{部署真正使用的}那条直线余量 $(b, v)$
(而不是先拟合误差曲线再转换,消除表示错配):行人/车辆余量全程比三代更瘦,
覆盖率合格,急刹 531 $\to$ 452(标定单项史上最大红利)。但旧决策层的参数是按
肥余量调的,直接换会小幅伤干净率——\textbf{裁决:不再给旧架构调参(补丁陷阱),
四代标定与新决策层配套一次调成}。

\subsection{待办清单(按优先级)}

\begin{enumerate}
  \item 新决策层的``刹停税''三修法,把到达率追平;
  \item 标准舒适指标接入每类目标的真实半径;
  \item 领域基线补齐(ORCA / DWA / MPC-CBF 对比);
  \item 统计封口:个别主张样本量不足($n{=}1$ 的姿态主张)、补``未触碰''确认折;
  \item 渲染面接入二/三代标定(消除两面不同步);
  \item 真引擎面执行器滞后走查、个别场景超时排查。
\end{enumerate}

% ============================================================
\section{论文定位(一段话)}

经双重文献查杀,两个卖点无近邻:
① \textbf{端到端组合定理}——``无分布假设的统计保证 $\times$ 连续时间严格证明
$\times$ 对实际执行路径 $\times$ 分类条件''四者的合取(单项都有人做,合取是空白);
② \textbf{规划器$\times$安全层的开放基准}——用严格连续时间证明作为与表示无关的
违约裁判。副卖点:``认证的优雅''(同等安全下更少急刹、更大反应裕度、姿态追平)。
按类别分层单独当主卖点会死(方法论已被 Mondrian CP 占据),降级为诚实审计段。
目标:RA-L 滚动投稿 + ICRA 2027。上游 EGO-Planner 的一个未定义行为 bug
已在移植中根修,可回报上游社区。

% ============================================================
\section{代码地图(去哪里找什么)}

仓库 \code{sando-core},唯一真相分支 \code{feat/bernstein-gate}
(\code{master} 上没有证明器和 EGO)。路径前缀 \code{M/} = \code{metaurban/}。

\begin{center}
\small
\begin{tabular}{@{}llp{7.2cm}@{}}
\toprule
职责 & 位置 & 一句话 \\
\midrule
感知前端 & \code{M/perception.py} & 视野锥/遮挡/漏检/噪声/无身份匹配 \\
跟踪预测 & \code{M/kf\_tracker.py} & 每目标一个匀加速卡尔曼滤波器 \\
误差标定 & \code{M/calibrate\_*.py}, \code{M/harvest\_v2.py} & 收割误差记录 $\to$ 三代标定文件 \\
安全证明器 & \code{cpp/include/sando\_cpp/bernstein\_cert.hpp} & 连续时间证明核(C++,7 微秒/次) \\
规划器桥 & \code{ego/} + \code{M/ego\_bridge.py} & 去 ROS 的 EGO-Planner + C 接口 \\
决策层 & \code{M/safety\_layer.py} & 候选锦标赛 + 粘性 + 紧急规避 \\
新决策层 & \code{M/local\_lattice.py} & ``认证集合内规划''(进行中) \\
回放评测 & \code{M/replay\_core.py}, \code{M/bench\_*.py} & 快速回放面 A/B \\
真引擎评测 & \code{M/ego\_mu\_bench.py}, \code{M/test\_ego\_mu.sh} & 版本阶梯回归(10 集/版本) \\
3D 渲染 & \code{M/render\_3d\_video.py} & 唯一合法出片产线 \\
场景工作台 & \code{M/scenario\_lib.py}, \code{M/populate.py} & 场景格式/生成/编辑/难度标定 \\
强化学习臂 & \code{M/train\_ppo\_*.py}, \code{M/shield.py} & 支撑验证(已归档) \\
权威文档 & \code{docs/safety-layer-spec.md} 等 & 规范/路线/定理/标定规范/问题清册 \\
\bottomrule
\end{tabular}
\end{center}

工程铁律(踩坑换来的,全文最后但每天都用):
标定必须匹配部署节奏;改跟踪/感知前后必跑字节级回归(\code{regress\_frozen\_ours.py});
两个 C 接口库不在 CMake 构建图里,改 C++ 后必须手动重编;
每一个算法改进主张都必须配 3D 对比视频取证;
汇报任何数字必须标明测量面;训练机 CPU 有硬件故障,长任务必须逐集落盘可续跑。

% ============================================================
\appendix

\section{术语对照表(本文用语 $\leftrightarrow$ 代码/文献名)}\label{app:glossary}

\begin{center}
\small
\begin{tabular}{@{}p{4.6cm}p{9.2cm}@{}}
\toprule
本文用语 & 代码 / 文献 / 记忆档案里的名字 \\
\midrule
预测误差余量 & conformal tube / 管子 / $T(\Delta)=q_0+v_{\text{eff}}\Delta$ \\
统计标定 & conformal calibration(split conformal) \\
一代标定 & \code{calib.json}(全局分位数) \\
二代标定 & FS3C-R / \code{calib\_v2.json}(分组细标) \\
三代标定(现役) & $\lambda$-SHAPE-H / \code{calib\_v3.json}(冻结形状+单缩放) \\
四代标定(待命) & LCP / v4b(直接优化部署直线) \\
安全证明 / 证明器 & Bernstein certificate / \code{bernstein\_cert.hpp} / cert \\
干净抵达 & clean arrival(到达 $\wedge$ 零碰撞 $\wedge$ 零紧急规避) \\
紧急规避 & evade / \code{evade\_setpoint} \\
候选锦标赛决策层 & \code{maneuver\_decide\_v2}(现役,即文中``旧决策层'') \\
新决策层 & CPL-v3 / \code{local\_lattice.py}(认证集合内规划) \\
现役最优栈 & V11(回放面)/ V7(真引擎面),一组开关的组合 \\
新目标限时 & \code{YOUNG\_TTL=2} \\
观测新鲜度修正 & \code{DELTA\_OVR=0.05} \\
爬行挡 & \code{SPEEDS\_CRAWL} \\
速度换证明窗 & \code{TAU\_SPEED}(信任窗=拍长+刹停时间) \\
最近点预测占据 & \code{CPA\_CLOUD}(CPA = 最近接近点) \\
逃逸压力提前触发 & \code{ESC\_TRIG} \\
刹车参考场 & \code{HOLD\_DECEL} \\
视野内全知对照 & cone oracle / 锥内全知 \\
$360^\circ$ 全知对照 & oracle \\
可交换 / 统计上同类 & exchangeability(conformal 保证的前提) \\
随时能刹停的退路性质 & 递归可行性(recursive feasibility) \\
测量面 & 快速回放面(replay)/ 真引擎面(MetaUrban)/ 3D 渲染面 \\
\bottomrule
\end{tabular}
\end{center}

\section{常用命令速查}

\begin{enumerate}
  \item \textbf{真引擎版本回归}(10 集 + 自动打印全版本进化表):\\
        \code{cd sando-core/metaurban \&\& ./test\_ego\_mu.sh <标签> [KEY=VAL ...]}
  \item \textbf{出对比视频}(唯一合法产线;三件套环境变量缺一段错误):\\
        \code{DISPLAY=:1 PYTHONPATH=<metaurban 仓> LD\_PRELOAD=<sando 环境 libstdc++>}\\
        \code{python render\_3d\_video.py --seed 7 --maneuver --clear\_spawn --mp4}\\
        (\code{--maneuver}=我们,\code{--ego}=原生基线;拼接用 \code{stitch\_ab.py})
  \item \textbf{三代标定}:\code{python calibrate\_v3.py}(产物 \code{calib\_v3.json})
  \item \textbf{C++ 重编}:\code{cmake --build build} 之外,两个 C 接口库必须手动
        \code{g++} 重编(见 \code{CLAUDE.md} 配方)
  \item \textbf{改动前后字节级回归}:\code{regress\_frozen\_ours.py}
\end{enumerate}

\end{document}
